% !TEX TS-program = pdflatex
% !TEX encoding = UTF-8 Unicode

% This file is a template using the "beamer" package to create slides for a talk or presentation
% - Giving a talk on some subject.
% - The talk is between 15min and 45min long.
% - Style is ornate.

% MODIFIED by Jonathan Kew, 2008-07-06
% The header comments and encoding in this file were modified for inclusion with TeXworks.
% The content is otherwise unchanged from the original distributed with the beamer package.

\documentclass{beamer}
\setbeamercovered{invisible}



% Copyright 2004 by Till Tantau <tantau@users.sourceforge.net>.
%
% In principle, this file can be redistributed and/or modified under
% the terms of the GNU Public License, version 2.
%
% However, this file is supposed to be a template to be modified
% for your own needs. For this reason, if you use this file as a
% template and not specifically distribute it as part of a another
% package/program, I grant the extra permission to freely copy and
% modify this file as you see fit and even to delete this copyright
% notice. 


\mode<presentation>
{
  \usetheme{Warsaw}
  % or ...

  % or whatever (possibly just delete it)
}


\usepackage[english]{babel}
% or whatever

\usepackage[utf8]{inputenc}
% or whatever

\usepackage{times}
\usepackage[T1]{fontenc}
% Or whatever. Note that the encoding and the font should match. If T1
% does not look nice, try deleting the line with the fontenc.

%%% MATH RELATED

\input{macros.tex}

\title{MATH 350-2 Advanced Calculus} 
\subtitle
{} % (optional)

\author[W.R. Casper] % (optional, use only with lots of authors)
{W.R. Casper}
% - Use the \inst{?} command only if the authors have different
%   affiliation.

\institute[California State University Fullerton] % (optional, but mostly needed)
{
  Department of Mathematics\\
  California State University Fullerton}
% - Use the \inst command only if there are several affiliations.
% - Keep it simple, no one is interested in your street address.

\subject{Talks}
% This is only inserted into the PDF information catalog. Can be left
% out. 



% If you have a file called "university-logo-filename.xxx", where xxx
% is a graphic format that can be processed by latex or pdflatex,
% resp., then you can add a logo as follows:

% \pgfdeclareimage[height=0.5cm]{university-logo}{university-logo-filename}
% \logo{\pgfuseimage{university-logo}}



% Delete this, if you do not want the table of contents to pop up at
% the beginning of each subsection:
\AtBeginSubsection[]
{
  \begin{frame}<beamer>{Outline}
    \tableofcontents[currentsection,currentsubsection]
  \end{frame}
}


% If you wish to uncover everything in a step-wise fashion, uncomment
% the following command: 

%\beamerdefaultoverlayspecification{<+->}


\begin{document}

\begin{frame}
  \titlepage
\end{frame}

\begin{frame}{Outline}
  \tableofcontents
  % You might wish to add the option [pausesections]
\end{frame}

% Since this a solution template for a generic talk, very little can
% be said about how it should be structured. However, the talk length
% of between 15min and 45min and the theme suggest that you stick to
% the following rules:  

% - Exactly two or three sections (other than the summary).
% - At *most* three subsections per section.
% - Talk about 30s to 2min per frame. So there should be between about
%   15 and 30 frames, all told.

\section{Real Analysis Lecture 14}

\subsection{Complete metric spaces}

\begin{frame}{Complete spaces}
\begin{thm}
Let $(M,d)$ be a metric space.
\pause
Then for any compact subset $S\subseteq M$, the subspace $(S,d)$ is complete.
\end{thm}
\pause
\begin{proof}
\pause
Let $(M,d)$ be a metric space and $S\subseteq M$ compact.\\
\pause
To show that $(S,d)$ is complete, we must show that every Cauchy sequence in $S$ converges to a value in $S$.\\
\pause
Suppose that $\{x_n\}$ is a Cauchy sequence in $(S,d)$.\\
\pause
If the set $X$ of all values of $\{x_n\}$ is finite, then the Cauchy condition implies $\{x_n\}$ converges.\\
\pause
Assume instead $X$ is infinite.\\
\pause
TIME OUT!
\end{proof}
\end{frame}

\begin{frame}{Complete spaces}
\begin{lem}
Let $(M,d)$ be a metric space and let $S\subseteq M$ be compact.
\pause
Then for any infinite subset $X\subseteq S$, the set $S$ has an accumulation point of $X$.
\end{lem}
\pause
\begin{proof}
\pause
To see this, suppose otherwise.\\
\pause
Then for all $x\in S$, there exists $r_x > 0$ such that $B_S(x,r_x)\cap X$ is empty if $x\notin X$ or equal to $\{x\}$ if $x\in X$.\\
\pause
The open cover $\{U_i: i\in I\}$ with $I=S$ and $U_i = B_S(i; r_i)$ of $S$ has no finite subcover.\\
\pause
This is a contradiction, so $S$ has an accumulation point of $X$.
\end{proof}
\end{frame}

\begin{frame}{Complete spaces}
\begin{thm}
Let $(M,d)$ be a metric space.
Then for any compact subset $S\subseteq M$, the subspace $(S,d)$ is complete.
\end{thm}
\begin{proof}
\pause
Let $L$ be an accumulation point of $X$ in $S$.\\
\pause
Let $\epsilon > 0$.\\
\pause
Then $\exists N>0$ such that $m,n \geq N$ implies $d(x_m,x_n) < \epsilon/2$.\\
\pause
Moreover, $B_S(L,\epsilon/2)$ contains infinitely many points of $X$.\\
\pause
Hence it contains $x_m$ for some $m\geq N$.\\
\pause
It follows that for all $n\geq N$,
$$d(x_n,L)\leq d(x_n,x_m) + d(x_m,L) < \epsilon/2 + \epsilon/2 = \epsilon.$$
\pause
Since $\epsilon > 0$ was arbitrary, this proves $\{x_n\}$ converges to $L$.
\end{proof}
\end{frame}

\subsection{Limit of a function}

\begin{frame}{Back in calculus}
The limit as $x\rightarrow a$ of $f(x)$ is $L$ means:\\
\pause
as $x$ gets closer and closer to $a$, 
\pause
$f(x)$ gets closer and closer to $L$.\\
\pause
\textbf{Rigorous:}
For all $\epsilon > 0$ there exists $\delta > 0$ such that
$$0 < |x-a| < \delta\ \ \Rightarrow\ \ |f(x)-L| < \epsilon.$$
\end{frame}


\begin{frame}{Limit of a function}
Let $(S,d_S)$ and $(T,d_T)$ be metric spaces and $A\subseteq S$.
\pause
Also, let $f: A\rightarrow T$ be a function.
\pause
\begin{defn}
\pause
Then for any accumulation point $p$ of $A$ and element $L\in T$, we say that the \textbf{limit as $x$ approaches $p$ of $f(x)$ is $L$}
\pause
 and write
$$\lim_{x\rightarrow p} f(x) = L$$
\pause
if for all $\epsilon > 0$ there exists $\delta > 0$ such that
$$0 < d_S(x,p) < \delta\ \Rightarrow\ d_T(f(x),L) < \epsilon.$$
\end{defn}
\end{frame}

\begin{frame}{Limit of a function}
Let $(S,d_S)$ and $(T,d_T)$ be metric spaces and $A\subseteq S$.
\pause
Also, let $f: A\rightarrow T$ be a function.
\pause
\begin{defn}
\pause
Then for any accumulation point $p$ of $A$ and element $L\in T$, we say that the \textbf{limit as $x$ approaches $p$ of $f(x)$ is $L$}
\pause
 and write
$$\lim_{x\rightarrow p} f(x) = L$$
\pause
if for all $\epsilon > 0$ there exists $\delta > 0$ such that
$$x\in B_S(p;\delta)\backslash\{p\}\ \Rightarrow\ f(x)\in B_T(L;\epsilon).$$
\end{defn}
\end{frame}

\begin{frame}{Limit of a function}
Let $(S,d_S)$ and $(T,d_T)$ be metric spaces and $A\subseteq S$.
\pause
Also, let $f: A\rightarrow T$ be a function.
\pause
\begin{defn}
\pause
Then for any accumulation point $p$ of $A$ and element $L\in T$, we say that the \textbf{limit as $x$ approaches $p$ of $f(x)$ is $L$}
\pause
 and write
$$\lim_{x\rightarrow p} f(x) = L$$
\pause
if for all $\epsilon > 0$ there exists $\delta > 0$ such that
$$f(B_S(p;\delta)\backslash\{p\})\subseteq B_T(L;\epsilon).$$
\end{defn}
\end{frame}



\end{document}


